% Load in the preamble; place the two commands at the indicated manuscript locations.
% \HSevenResults: after the Figure 5 structural-example subsection.
% \HSevenMethods: immediately before Audit and reproducibility.
\newcommand{\HSevenResults}{%
\subsection{Local structural context adds little transferable rank information}

The structural examples motivate a further question: can local atomic information
explain differences in response beyond those captured by global pore descriptors?
We addressed this after the pair catalogue was frozen, using the same high-minus-low
orientation and primary CH$_4$ endpoint. Lightweight CrystalNN-derived descriptors
summarized metal coordination, neighbour chemistry, bond-length variation, local
order parameters and atomic-degree statistics. Fingerprinting succeeded for 22,137
of 22,144 structures, leaving 11,065 pairs for an identical comparison of global
geometry, local descriptors and their combination (\cref{fig:local-context}A--B).
These models predict the within-pair uptake difference; they neither redefine the
matching nor replace the primary population estimator.

Under random five-fold cross-validation, mean Spearman correlation increased from
0.8639 for geometry alone to 0.8666 for the combined representation, while mean
absolute error decreased from \SI{0.6724}{\milli\mole\per\gram} to
\SI{0.6453}{\milli\mole\per\gram}, a 4.0\% reduction. Both measures improved
in each random fold. The local-only model was less predictive, with mean Spearman
correlation 0.6543. The added rank information was therefore small relative to that
already available from the eight pore descriptors and density.

Holding out entire topology groups changed the comparison. Mean Spearman
correlation was 0.7490 for geometry alone, 0.3257 for local descriptors, and 0.7372
for their combination. The combined model had lower correlation in all five
topology-held-out folds (\cref{fig:local-context}D). Mean absolute error changed
only slightly, from \SI{0.8357}{\milli\mole\per\gram} to
\SI{0.8330}{\milli\mole\per\gram}. Thus, the small random-CV improvement did
not translate into improved rank prediction across held-out topologies.

\mainfigure{Figure_6.pdf}{fig:local-context}
  {Incremental predictive information from local structural context}
  {(A) Frozen-pair feature differences; (B) geometry-only, local-only and combined
   models; (C) held-out feature attribution; (D) within-fold incremental rank performance.}
  {\HSevenCaption}

Held-out SHAP attribution was dominated by density and pore descriptors
(\cref{fig:local-context}C). Raw metal count was the highest-ranked local input,
but its dependence on unit-cell representation prevents interpreting that
attribution as a coordination mechanism. Moreover, tetrahedral and square-planar
feature differences were unavailable for 85.5\% of usable pairs, and octahedral
differences for 93.0\%; these values were imputed within training folds. The
implemented metal--metal distance summary also excludes periodic copies of the
same metal site. These representation limitations bound the conclusion: the
implemented local fingerprint provides little additional transferable rank
information for this endpoint, but the analysis does not establish that local
coordination in general is uninformative. SHAP describes model use rather than
causal adsorption mechanisms.


}

\newcommand{\HSevenMethods}{%
\subsection{Post-freeze local structural descriptors and model comparison}

The local-context analysis retained all 11,072 frozen pairs, the 10\% eligibility
criterion, pair orientation, response archetypes and primary endpoint. Pymatgen's
CrystalNN neighbour finder was applied to each required CIF with unweighted
coordination numbers, \texttt{cation\_anion=False}, distance cutoffs (0.5, 1.0),
electronegativity weight 1.5, porous adjustment enabled, and a
\SI{4.5}{\angstrom} search cutoff. No oxidation states were assigned. The
19 local inputs comprised metal count and fraction; mean and standard deviation
of metal coordination number and metal--neighbour distance; mean per-metal
bond-length coefficient of variation; O-, N- and S-neighbour fractions; mean
absolute metal--neighbour electronegativity difference; mean tetrahedral,
square-planar and octahedral order parameters where returned; atomic-degree mean,
standard deviation and coefficient of variation; and mean and standard deviation
of nearest distances between distinct metal-site indices. Bond-length variation
was not treated as a complete measure of angular distortion.

Each descriptor was converted to a high-minus-low difference. Three fixed LightGBM
regressors used nine global inputs (eight pore descriptors plus density), 19 local
inputs, or their 28-input union. The models used an absolute-error objective,
400 trees, learning rate 0.035, at most 24 leaves and depth 6, a minimum of 40
samples per leaf, feature-sampling fraction 0.85, L1/L2 penalties 0.10/0.50, and
seed 17, without hyperparameter search or outcome-based neighbour tuning. All
models used the same 11,065 pairs with successful fingerprints for both members.
A predefined minimum coverage of 90\% at both framework and pair levels was
exceeded; seven structures failed because Voronoi neighbours could not be found.
No pairs were rematched or replaced.

Evaluation used shuffled five-fold cross-validation (seed 17) and five-fold
GroupKFold with recorded topology as the grouping variable. No framework appears
in more than one frozen pair. Grouped test folds contained 4,659, 4,219, 1,051,
653 and 483 pairs, respectively; the first four corresponded to pcu, fsc, nbo and
sra, while the last combined the remaining 45 topologies. Missing feature values
were imputed by training-fold medians, with zero used only when an entire training
column was missing. Spearman correlation, mean absolute error and $R^2$ were
computed on held-out pairs and summarized by arithmetic means and standard
deviations across folds. TreeSHAP contributions were computed on held-out pairs
using LightGBM's contribution output. The supplementary saved pair-bootstrap
distribution uses 2,000 resamples of pooled out-of-fold predictions; it is a
different statistic from mean fold Spearman and does not quantify uncertainty
over newly sampled topology groups.

}
